%%%%%%%%%%%%%%%%%%%%%%%%%%%%%%%%%%%%%%%%%%%%%%%%%%%%%%%%%%%%%%%%%%%%%%%%%%%%%%%
%  s02-background.tex —— 第 2 节 辛拓扑背景（对应原文 pp.3--6；源码第 394--507 行）
%  原文：Ivan Smith, “A Symplectic Prolegomenon”, arXiv:1401.0269v2, Section 2
%%%%%%%%%%%%%%%%%%%%%%%%%%%%%%%%%%%%%%%%%%%%%%%%%%%%%%%%%%%%%%%%%%%%%%%%%%%%%%%

\section{辛拓扑背景}\label{Sec:Symplectic}

\begin{keypoints}
\begin{itemize}
  \item \textbf{核心图像}：辛流形是带有闭的、非退化 2-形式 $\omega$ 的偶维流形。$\omega$ 的非退化性
        给出等同 $T^*M \cong TM$，使每个光滑函数 $H$ 都生成保持 $\omega$ 的 Hamilton 流；
        Darboux--Moser 定理说明辛几何没有局部不变量，因此有意义的拓扑问题都是整体问题。
  \item \textbf{基本工具}：驯服 $\omega$ 的近复结构构成一个\emph{可缩}空间 $\scrJ$；对一般的 $J$，
        全纯曲线的模空间是有限维光滑流形，并容许 Gromov 紧化。Chern 类与 Maslov 类通过指标理论
        给出这些模空间的（虚拟）维数。
  \item \textbf{动力系统与相交}：Hamilton 方程的解恰好是余切纤维在时间 $1$ 流下的像与另一条余切纤维
        的交点；而辛同胚的不动点等同于乘积辛流形中两个拉格朗日子流形的交点。
  \item \textbf{与其他领域的联系}：辛几何是少数几种“自然的”无穷维几何之一；代数族的单值通过
        Dehn 扭转表现为辛现象；Donaldson 关于 4-流形的猜想把辛拓扑与光滑拓扑的异常特征联系起来。
  \item \textbf{镜像对称}：同调镜像对称猜想 $D^{\pi}\scrF(X) \simeq D^b(\Cech{X})$ 断言辛侧
        导出 Fukaya 范畴与复侧凝聚层的导出范畴等价，这是全文的主要动机。
\end{itemize}
\end{keypoints}

\subsection{基本概念}

辛流形 $(M^{2n},\omega)$ 是偶维流形，它可以由 $\bR^{2n}$ 中的开集通过导数落在子群 $Sp_{2n}(\bR) = \left\{A \ | \ A^t J A = J \right\} \leq GL_{2n}(\bR)$ 中的微分同胚粘合而成，其中 $J = \left(\begin{array}{cc} 0 & -1 \\ 1 & 0 \end{array} \right)$ 是欧氏空间上固定的斜对称形式。在 $\bC^n$ 上取坐标 $z_j = x_j + iy_j$，$1\leq j \leq n$，并令 $\omega_{\st} = \sum_j dx_j \wedge dy_j$：这些微分同胚正是保持 2-形式 $\omega_{\st}$ 的微分同胚；而事实上，一个流形是辛的当且仅当它容许一个闭的（$d\omega = 0$）、非退化的（$\omega^n > 0$）2-形式 $\omega$。仅此一点就已经排除掉了一些拓扑，首先便是球面 $S^4$。

辛几何最重要的单一特征是：经由 $\omega$ 的（非退化性所给出的）等同 $T^*M \cong TM$，任意光滑函数 $H: M \rightarrow \bR$ 都定义了一个向量场 $X_H = \iota_{\omega}(dH)$，其流 $\phi^t_H$ 保持辛结构。于是辛流形具有无穷维的局部可迁对称群，并且不存在类似 Riemann 曲率的局部不变量。Darboux--Moser 定理强化了这一点，它断言：(i) 任意辛 $2n$-流形局部同构于 $(B^{2n}, \omega_{\st})$；(ii) 保持上同调类的 $\omega$ 的形变可以提升为整体辛同胚，至少当 $M$ 紧时如此。因此，辛几何中有意思的拓扑问题本质上是整体性的。换一种世界观看，也可以说：任何辛流形局部都与欧氏空间一样复杂，而关于欧氏空间，仍有大量问题悬而未决。

% [Tutor Note: 上面 Darboux--Moser 定理的 (i)(ii) 两条都是 Moser 同伦论证的推论：若闭形式 $\omega_t$
%   保持上同调类，可写 $\omega_t=\omega_0+d\alpha_t$，并由 $X_t=\iota_{\omega_t}^{-1}(\dot\alpha_t)$
%   积分得到同痕。另请注意“非退化”只给出 $T^*M\cong TM$ 这一等同，局部标准形来自 Darboux 定理，
%   二者不可混为一谈。]

两类特别重要的辛流形是：
\begin{enumerate}
\item （准）射影簇 $X \subset \bP^N$，配以 Fubini--Study 形式 $\omega_{\operatorname{FS}}$ 的限制；特别地，任意有限型的定向二维曲面都在此列，此时辛形式不过是一种度量面积的方式；
\item 余切丛 $T^*Q$，配以典范形式 $\omega =d\theta$，其中 $\theta$ 是 Liouville 1-形式；更一般地，还有 Stein 流形，即 $\bC^N$ 中闭的、真嵌入的全纯子流形，等价地说，是配有穷竭的严格多重次调和函数 $h$ 的复流形，此时 $\omega = -dd^ch$。
% \item 李群的余伴随轨道，配以 Kostant--Kirillov--Souriau 辛形式。
\end{enumerate}

% [Tutor Note: 这三个包含关系里，第一条来自 Kähler 形式本身就是辛形式；第二条是定义使然——
%   辛流形按定义就是光滑流形（驯服近复结构还额外给出几乎复结构）。$4$ 维时 Gompf 的手术构造
%   给出可证明非代数的辛流形，说明第一个包含是严格的；而“辛 vs 光滑”的鸿沟需要规范理论等工具。]

辛流形足够灵活，可以容许“剪切-粘贴”型的手术 \cite{Gompf}；其中许多四维的例子可以证明不是代数的；辛拓扑中的一个主导主题，就是理解包含关系 $\{$K\"ahler$\} \subset \{$Symplectic$\} \subset \{$Smooth$\}$。对于开流形，辛结构有良好的存在性理论，但唯一性问题依然微妙：欧氏空间上存在可数多个辛上互不相同的（有限型）Stein 结构，而识别标准结构在算法上是不可判定的 \cite{McLean:Computability}。

\subsection{进一步的概念}

% [Tutor Note: 可缩性在这里是实质的：不变量依赖辅助选取的 $J$，而空间可缩保证任意两个选取之间
%   的道路在同伦意义下唯一，于是带符号计数与分次不随 $J$ 改变。若换成一般的近复结构空间
%   （非可缩），这套论证立刻失效。]

辛几何的第二个本质特征是：任意辛流形都容许一个\emph{可缩}的空间 $\scrJ$，由驯服辛形式的近复结构组成；所谓驯服，是指对任意非零切向量 $v$ 都有 $\omega(v, Jv) > 0$。遵循一种熟悉的模式，相对于从可缩空间中所作的辅助选择而定义的离散（例如计数型）不变量，反映的是底下辛结构的内蕴性质。

% [Tutor Note: 排除“多重覆盖”曲线是为了横截性：多重覆盖以非正则的方式出现，会使模空间不再是
%   光滑流形、计数没有良好定义；处理它们需要虚拟扰动技术（见第 3 节）。]

子流形 $L\subset M$ 称为拉格朗日的，如果它是迷向的（$\omega|_L = 0$）且是半维的（即迷向子流形所可能具有的最大维数）。在上述几类辛流形中，拉格朗日子流形包括：(1) 定义在 $\bR$ 上的代数簇的实点轨迹；(2) $T^*Q$ 中的零截面与余切纤维（分别对应）。Gromov 的基本洞见是：对于一般的驯服近复结构，$M$ 中非多重覆盖的全纯曲线空间（或许带有拉格朗日边界条件）是有限维光滑流形，而且它们还容许由结点曲线给出的、具有几何意义的紧化。正是由此得到的不变量——即计数“伪全纯”曲线（从 Riemann 面出发、对某个驯服但未必可积的 $J$ 为 $J$-复线性的映射）的不变量——支配着这个领域。

% [Tutor Note: Maslov 类的符号与归一化约定极易出错：不同文献可能把 $\mu_L$ 取为 $2c_1$ 或
%   $c_1$、或相差一个整体符号；本段以 $H^2(M,L;\bZ)\to H^2(M;\bZ)$ 下 $\mu_L\mapsto 2c_1(M)$
%   为准。核对时请回到“沿圆盘边界拉格朗日切平面的绕数”这一几何定义，勿只比较代数公式。]

辛流形与拉格朗日子流形的离散不变量，包括一些特征类；通过指标理论，它们决定了全纯曲线模空间（在非空时）的虚拟维数，即期望维数。任意驯服近复结构 $J$ 的选择都使 $TM$ 成为复丛，因而有 Chern 类 $c_i(M) \in H^{2i}(M;\bZ)$。由于 $\bC^n$ 的拉格朗日子空间构成的 Grassmann 流形 $U(n) / O(n)$ 的基本群为 $\bZ$，任意拉格朗日子流形都有一个 Maslov 类 $\mu_L: \pi_2(M,L) \rightarrow \bZ$，它度量拉格朗日切平面沿圆盘 $D$ 的边界、相对于 $TM|_D$ 的某个平凡化所作的扭转；在映射 $H^2(M,L;\bZ) \rightarrow H^2(M;\bZ)$ 下，$\mu_L$ 映到 $2c_1(M)$。表示同调类 $A$ 的\emph{参数化}有理曲线 $\bP^1 \rightarrow M^{2n}$ 的模空间具有实维数 $2\langle c_1(M),A\rangle + 2n$，于是除以 $\bP SL_2(\bC)$ 后，非参数化曲线的空间维数为 $2\langle c_1(M), A\rangle + 2n - 6$。$L$ 上以 $L$ 为边界、相对同调类为 $u \in \pi_2(M,L)$ 的参数化全纯圆盘构成的空间，具有虚拟实维数 $n + \mu(u)$，并带有 3 维实群 $\bP SL_2(\bR)$ 的作用。

\subsection{语境}

究竟为什么要研究辛拓扑？

\subsubsection{李理论}

除了众所周知的有限维李群分类之外，Lie 与 Cartan 还考虑过欧氏空间的无穷维对称群。为消除冗余，他们考虑在 $\bR^n$ 上局部可迁地作用、且不保持任何由平行平面 $\bR^k \subset \bR^n$ 给出的非平凡叶状结构的对称群。结果这样的伪群只有寥寥几种（现代论述见 \cite{SingerSternberg}）：全体微分同胚群、保体积微分同胚群及其共形类似物、偶数维中的辛变换或共形辛变换、奇数维中的接触变换，以及这些群的全纯类似物。于是，辛几何成为少数几种“自然的”无穷维几何之一（与有限维的分类相对照：$\bR^n$ 的对称群没有“例外”的无穷维情形）。

\subsubsection{动力学}

固定光滑函数 $H: \bR^{2n} \rightarrow \bR$。Hamilton 方程是
\[
\frac{\partial H}{\partial q} = -\dot{p}; \ \frac{\partial H}{\partial p} = \dot{q}
\]
其中 $(q,p) \in \bR^n \times \bR^n = T^*\bR^n$。若 $H = p^2/2 + U(q)$，而 $U: \bR^n \rightarrow \bR$ 是势函数，则这些方程支配一个质量为 1、位置-动量坐标为 $(q,p) \in T^*\bR^n$、受外力 $ -\partial U /\partial q$ 作用的质点的演化。恒等式
$
\frac{dH}{dt} =  \frac{\partial H}{\partial q} \frac{dq}{dt} + \frac{\partial H}{\partial p} \frac{dp}{dt} = 0
$
给出能量 $H$ 的守恒；而经典 Poincar\'e 回复现象则源自如下事实：系统的演化保持体积，即保持测度 $\omega_{\st}^n = \prod_j dp_j dq_j$；但事实上，演化保持的是底下的辛形式 $\omega_{\st}$。因此，经典 Hamilton 动力学无非是沿着 $\Symp(\bR^{2n})$ 中的一条道路行进。

% [Tutor Note: $\phi_H^1(T_x^*)$ 指余切纤维 $T_x^*Q$ 在时间 $1$ Hamilton 流下的像；$x$ 到 $y$
%   的轨迹恰是它与 $T_y^*$ 的交点。这一“动力学 = 拉格朗日相交”的翻译，正是 Floer 理论把
%   Hamilton 轨迹问题化为拉格朗日相交问题的出发点。]

给定两点 $x,y \in Q$ 以及时间依赖的 Hamilton 函数 $H_t: T^*Q \rightarrow \bR$，拉格朗日子流形 $\phi_H^1(T_x^*)$ 与 $T_y^*$ 的交点恰好对应该系统从 $x$ 到 $y$ 的轨迹；因此从动力学观点看，拉格朗日相交从一开始就是自然的。

\subsubsection{代数族的几何\label{Sec:AlgFamily}}

设 $f: X \rightarrow B$ 是射影簇之间的全纯映射，在 $B^0 \subset B$ 上光滑。则 $f$ 在 $B^0$ 上定义局部平凡的纤维化。多数情形下，$f$ 的光滑纤维作为复流形互不相同，但 Moser 定理蕴涵它们作为辛流形同构；纤维化的结构群约化为 $\Symp(f^{-1}(b))$。例如考虑二次映射 $\pi: \bC^3 \rightarrow \bC$，$(z_1,z_2,z_3) \mapsto \sum z_j^2$，它在 $\bC^*$ 上光滑。远离原点的纤维是光滑的仿射二次曲线，辛同胚于 $(T^*S^2, d\theta)$。该纤维化的单值是关于拉格朗日消没圈 $S^2 \subset T^*S^2$ 的一个 \emph{Dehn 扭转}。一个统一了 Arnol'd、Kronheimer 与 Seidel 洞见的基本观察\footnote{$\Diff_{ct}$ 与 $\Symp_{ct}$ 分别表示紧支微分同胚群与紧支辛同胚群。}是：这个单值在 $\pi_0 \Diff_{ct}(T^*S^2)$ 中是 2 阶的，但在 $\pi_0 \Symp_{ct}(T^*S^2)$ 中是无限阶的 \cite{Seidel:twist}，见第 \ref{Sec:Appl1} 节。因而，从辛单值过渡到光滑单值会丢失信息。这一现象无处不在；在代数几何中涉及模的那些部分，单值的辛方面似乎是不可避免的。

辛同胚 $\phi: X \rightarrow X$ 的不动点，恰好是乘积 $(X\times X, \omega \oplus -\omega)$ 中拉格朗日子流形 $\Delta_X, \Gamma_{\phi}$ 的交点，这把辛单值的问题与拉格朗日相交联系了起来。

\subsubsection{低维拓扑}

低维拓扑中的许多不变量都有辛的解释（Alexander 多项式、Seiberg--Witten 不变量等）。Donaldson 猜想 \cite[p.437]{McD-S}：两个同胚的光滑辛 4-流形 $X,Y$ 微分同胚，当且仅当 $X\times S^2$ 与 $Y\times S^2$ 上的乘积辛结构形变等价（即属于辛结构空间的同一连通分支）；这暗示辛拓扑与 4-流形拓扑的“奇异”特征从根本上纠缠在一起。

\subsection{若干未解问题}

辛拓扑中的许多基本问题，迄今似乎完全无从下手。例如，当 $n>2$ 时，$\Symp_{ct}(\bR^{2n}, \omega_{\st})$ 是否可缩、是否连通、甚至是否只有可数多个连通分支，都属未知。迄今也不知道是否存在这样的例子：一个维数至少为 6 的单连通闭流形，它容许近复结构，并且拥有一个最高次幂非零的二次上同调类，却不带辛结构。我们并不否认这类问题的中心地位，但下文将集中于那些至少已取得某些进展的问题。

\subsubsection{拉格朗日子流形的拓扑}

这里的问题大致分为两类：对于“就在眼前”的拉格朗日子流形，我们能说什么（唯一性、可移开性……）；对于“看不见”的拉格朗日子流形，我们又能说什么（存在性、Maslov 类……）？

\begin{Question} $(\bC^n, \omega_{\st})$ 或 $(\bP^n, \omega_{\operatorname{FS}})$ 中的拉格朗日子流形有哪些？ \end{Question}

更具体地说，我们知道能拉格朗日嵌入 $(\bC^3,\omega_{\st})$ 的素 3-流形只有乘积 $S^1\times \Sigma_g$ \cite{Fukaya:LagC3}；对任意可定向 3-流形 $Y$，$Y\#(S^1\times S^2)$ 都可拉格朗日嵌入 $(\bC^3,\omega_{\st})$ \cite{YETI}；但对于哪些双曲 3-流形的连通和容许这样的拉格朗日嵌入，人们实质上还一无所知。

\begin{Question}\label{Q:Arnold} 若 $T^*L \cong_{\omega} T^*L'$ 是辛同胚的，那么 $L$ 与 $L'$ 是否必定微分同胚？ \end{Question}

Arnol'd 的“近旁拉格朗日子流形”猜想将预言“是”。Abouzaid 与 Kragh \cite{Kragh,Abouzaid:Maslov0} 已经证明 $L$ 与 $L'$ 必定同伦等价；在少数情形下——$L=S^{4n+1}$ 或 $L= (S^1\times S^{8n-1})$，见 \cite{Abouzaid:exotic,EkholmSmith}——人们知道 $T^*L$ 记得 $L$ 上光滑结构的某些侧面，即对某些同伦球面 $\Sigma$ 有 $T^*L \not\cong_{\omega} T^*(L\#\Sigma)$。

\begin{Question} 一个闭的 $(X,\omega)$，若 $\dim_{\bR}(X)>2$，能否包含无穷多个两两不相交的拉格朗日球面？ \end{Question}

对射影空间中的完全交，Hodge 理论给出关于\emph{结点}个数（从而也是同时消没圈个数）的界；但拉格朗日球面与消没圈之间的关系略显神秘，参见 \cite{CortiSmith}。

\subsubsection{映射类群与动力学}

辛映射类群 $\pi_0 \Symp(X)$ 是相当微妙的不变量：当 $\omega$ 的扰动改变其上同调类时（从而不受 Moser 定理控制），它可能发生剧烈变化。

\begin{Question}\label{Question:Monodromy} 设 $X_d \subset \bP^N$ 是次数为 $d$ 的超曲面，$d>2$，而 $\scrX_d$ 是所有这些超曲面的粗模空间；描述自然表示 $\pi_1(\scrX_d) \rightarrow \pi_0\Symp(X_d)$ 的（余）核。 \end{Question}

已知该映射的像无限；在这一般性下，它似乎不大可能是单射。余核就更加神秘了：$\pi_0\Symp(X_d)$ 是否有限生成，尚属未知。更靠近动力学一侧，人们可以寻找 Rivin 结果 \cite{Rivin} 的类似物：经典映射类群上的随机游走几乎必然给出伪 Anosov 元素。

\begin{Question} \label{Question:RandomWalk} 对 $\pi_0\Symp(X_d)$（$d>2$）上的随机游走，最终映射类的任意代表元的周期点个数是否以概率一指数增长？ \end{Question}

还有一圈相当不同的动力学问题，来自对\emph{通量}的探讨。一个元素 $\phi \in \Symp_0(M)$，即与恒等同痕的辛同胚，若它是某个 Hamilton 流 $\phi_H^t$ 的时间 1 映射，则称为 Hamilton 的。固定 $a\in H^1(M;\bR)$ 以及闭 1-形式的一条道路 $a_t$，满足 $\int_0^1 a_t dt = a$。这些 $a_t$ 对偶于辛向量场 $X_t$，后者生成同痕 $\phi^t$；而 $\phi^1 = \phi \in \Symp_0(M)$ 在 Hamilton 同痕意义下只依赖于 $a$，不依赖于道路 $(a_t)$。\emph{通量群} $\Gamma \subset H^1(M;\bR)$ 由那些使 $\phi^1 \simeq \id$（即与恒等 Hamilton 同痕）的 $a$ 组成。Ono 的一个深刻定理 \cite{Ono} 蕴涵 $\Gamma$ 是离散的。

\begin{Question} 哪些 $a\in H^1(M,\bR)$ 可由无不动点的辛同胚实现？ \end{Question}

% [Tutor Note: 这里依赖 Arnold 猜想（Hamilton 微分同胚的不动点下界，非退化情形下形如
%   $\#\mathrm{Fix}(\phi)\ge\sum_j\dim H^j(M;\bQ)$）的解决，故通量元素不能由无不动点的
%   辛同胚实现。注意 flux 探测 $\Symp_0(M)$ 的同痕信息，而 $\pi_0\Symp(M)$ 是整体映射类群。]

由（Arnold 猜想的解决）\cite{FOno, LiuTian} 可知，$\Gamma$ 中的元素不可能由无不动点的辛同胚实现。注意，通量探测的是 $\Symp_0(M)$，而不是 $\pi_0 \Symp(M)$，这与问题 \ref{Question:RandomWalk} 形成对照。

\subsection{镜像对称的现象与哲学}\label{Sec:Mirror}

% [Tutor Note: “无障碍”（unobstructed）指 $m_0=0$，或选定界定上链后 $m_0$ 可被消去；只有这样的
%   拉格朗日子流形才是 Fukaya 范畴的对象，一般情形下得到的是 $A_\infty$-模。所谓形式代数扩充，
%   即取三角闭包并作幂等完备化（split-closure）。]

在其最初的表述中，Kontsevich 的同调镜像对称猜想 \cite{Kontsevich} 断言
\begin{equation} \label{Eqn:HMS}
D^{\pi}\scrF(X) \simeq D^b(\Cech{X})
\end{equation}
其中 $X, \Cech{X}$ 是一对互为镜像的 Calabi--Yau 流形，$D^b(\Cech{X})$ 表示有界凝聚层导出范畴，而 $D^{\pi}\scrF(X)$ 是分裂闭包的导出 Fukaya 范畴。前者以代数向量丛的有限长复形 $\mathcal{E}^{\bullet} \rightarrow \Cech{X}$ 为对象，态射为 Cech 复形 $\Cech{C}^*(\EuScript{U}; \mathcal{H}om(\mathcal{E}^{\bullet}, \mathcal{F}^{\bullet}))$，其中 $\EuScript{U}$ 是 $\Cech{X}$ 的一个固定仿射开覆盖；后者则由一个形式代数的扩充过程得到，其对象是 $X$ 的“无障碍”拉格朗日子流形，态射是下文将引入的 Floer 上同调群，见第 \ref{Sec:Floer} 节与第 \ref{Sec:Fukaya} 节。这一表述掩盖了一个事实：一般而言，给定 $X$ 并没有明显的方式产生 $\Cech{X}$，而且它可能并不唯一。（Gross 与 Siebert \cite{GrossSiebert} 给出的更好表述，涉及某类 Calabi--Yau 环簇退化上的一个对合对称性；这些退化的中心纤维应当是以直接组合的方式互为对偶的环簇之并。）

% [Tutor Note: 方向约定务必留意：左边是辛侧 $X$ 的导出 Fukaya 范畴，右边是复几何侧镜像
%   $\Cech{X}$ 的凝聚层导出范畴；不少作者把两侧对调写成 $D^b(X)\simeq D^{\pi}\scrF(\Cech{X})$。
%   上标 $\pi$ 表示 split-closure（幂等完备化），不是 “projective”；去掉它会使命题不成立。]

% [Tutor Note: 此处的“Legendre 变换”是整仿射流形上的组合对偶操作，与经典力学中函数到函数的
%   Legendre 变换只是形式类比；它实现的是 SYZ 意义下的 T-对偶，即环面纤维与其对偶环面互换。]

几何上，按照 Strominger--Yau--Zaslow \cite{SYZ} 的看法，$X$ 与 $\Cech{X}$ 之间的猜想关系是：它们带有相互对偶的拉格朗日环面纤维化（允许奇异纤维）；这种拉格朗日纤维化的底空间是带奇点的整仿射流形，而环簇对偶来自这类流形上的“Legendre 变换”。由此观点，镜像对称已被推广到远为广泛的语境：给定一个带有效反典范除子 $D\subset X$ 的簇 $X$，人们在 $U=X\backslash D$ 上寻找拉格朗日纤维化，并取其对偶 $\Cech{U}$。从 $U$ 回到 $X$ 的紧化，镜像地对应于开启一个全纯“势”函数 $W: \Cech{U}\rightarrow \bC$——它由以拉格朗日环面纤维为边界的全纯圆盘计数构成——反之亦然。（在原来的 Calabi--Yau 情形，$D = \emptyset$，$W$ 为常数。）镜子的两侧仍然有可以相互比较的拉格朗日范畴与层范畴：例如，若 $X$ 是 Fano 的，人们会问是否有 $D^{\pi}\scrF(X) \simeq D^b_{sing}(W:\Cech{U} \rightarrow \bC)$，右端是 $W$ 的矩阵分解的同伦范畴 \cite{Orlov:Dbsing}。

阐明这一点就离题太远了：为着提供动机，只需知道，存在若干代数几何的途径来计算某些范畴，这些范畴被猜想——有时也被证明——等价于由紧拉格朗日子流形构造出的范畴，尽管人们并无任何现实的希望去分类甚至枚举这些拉格朗日子流形。反过来，代数几何中感兴趣的范畴也可以用辛方法来研究：$X$ 的辛自同构作用在 $D^{\pi}\scrF(X)\simeq D^b(\Cech{X})$ 上，而 $X$ 的辛映射类群往往比 $\Cech{X}$ 的代数自同构群丰富得多。

\begin{Notes} 关于辛拓扑的基础，见 \cite{ACdaS, McD-S}；关于动力学与拓扑，见 \cite{HZ, EOSS}；关于镜像对称的概览，见 \cite{AspinwallEtAl, Ballard, Gross:tropical}。\end{Notes}
